\thispagestyle{plain}
\phantomsection
\addcontentsline{toc}{chapter}{中文摘要}
\centerline{\zihao{-3}\heiti 中文摘要}

\linespread{1.4}\zihao{-4}\bigskip

本文围绕多测站空间目标定位问题，研究了融合测角与测距观测的非线性参数估计方法。针对观测量类型异构、噪声尺度不一致、角度残差具有周期性等特点，建立了统一的加权最小二乘目标函数，并基于目标函数的一阶与二阶信息给出了两类求解框架：牛顿迭代法与法方程迭代法。为提升迭代稳定性，本文进一步采用三测站测距几何交会方法构造初始值，使优化过程具备可复现性和较强收敛性。

在模型层面，本文给出了斜距函数、方位角函数、俯仰角函数的解析梯度与海森矩阵表达，并讨论了方位角差值的 $2\pi$ 周期修正策略。算法层面，牛顿法用于直接求解非线性最优化问题，法方程迭代法用于构造线性化增量系统，二者在程序实现中共享观测模型与导数模块。实验层面，依据课程代码设置了 4 组测角观测和 4 组测距观测，对定位精度、迭代收敛、矩阵条件性、误差传播、置信椭球和鲁棒性表现进行了系统分析。

结果表明：两种方法均能稳定收敛到同一最优位置估计；在当前几何构型下，水平分量估计精度高于垂直分量；法化矩阵逆矩阵与牛顿法收敛点海森矩阵逆具有一致量级，可用于协方差近似；通过卡方分位数和谱分解构造的置信椭球能够直观反映估计不确定性的方向性。进一步地，本文通过多组噪声水平与站位扰动实验验证了模型的适用范围，并给出了工程实现中的参数设置建议与扩展方向。

本文完成了“建模--推导--算法--实现--评估”完整闭环，形成了一套面向课程项目与工程原型的空间定位求解流程。研究成果可为雷达光电协同观测、无人机目标定位和多传感器信息融合场景提供参考。

\bigskip
\noindent{\zihao{4}\heiti 关键词：}
空间目标定位；加权最小二乘；牛顿法；法方程；置信椭球；误差分析
